\section{Appendix D: Proof assembly and verification note}

The preceding auxiliary results supply the regularity, assignment, and converse ingredients of the capacity characterization. This section collects their implications for the minimax criterion, the dimension threshold, and causal precision, as stated in \cref{obj:thm:capacity-answer,obj:thm:log-free-bivariate-capacity}. The following complete characterization assembles the common assignment procedure, the independent-prior certificate, and the exact causal interpretation.

\begin{theoremv}{thm:capacity-answer}[Complete capacity and excess characterization]*
For every pair of integers \(n,d\geq2\), let \(B\), the coordinate-pair count, and \(a_s(n,d)\), the comparison scale, be
\[
B=\binom d2,\qquad
a_s(n,d)=\min\{1,(B/n)^s\}.
\]
There is a single fair, outcome-independent Borel assignment kernel \(\pi^\star_{n,d}\) in the design class of \cref{obj:def:design-class}, using the original \(n\) units and independent of \(s\). For \(n\leq B\), this kernel assigns independent fair signs. For \(B<n\), it is the law of the prescribed spectral signing procedure, using independent reflection coins, an independent common Haar shift, and independent uniform rounding seeds. The rounding procedure starts at the zero fractional assignment and uses the first \(K=\lfloor n/4\rfloor\) prescribed real spectral rows evaluated at the reflected and commonly shifted sample. For every realization of the reflection coins and common shift, and every sequence of rounding seeds in \([0,1]\), its fractional assignment after \(n\) iterations equals the resulting assignment sign vector:
\[
u_i=Z_i,\qquad i=1,\ldots,n.
\]

For every \(0<s\leq1\), the lower comparison scale satisfies \(b_s(n,d)=a_s(n,d)\). With
\[
c_s=\frac{2}{(48+32\pi^2)16384^s},
\qquad
c_1=\frac{2}{(48+32\pi^2)16384},
\]
the capacity and signing loss of \cref{obj:def:criterion}, over the function class of \cref{obj:def:sobolev-class}, satisfy
\[
0<c_1\leq c_s,\qquad
c_sa_s(n,d)
\leq R_s(n,d)
\leq
\sup_{m\in\mathcal H_{d,s}}
\mathcal L_{n,d,s}(\pi^\star_{n,d},m)
\leq3072a_s(n,d).
\]
These inequalities hold simultaneously for the same kernel \(\pi^\star_{n,d}\).

Every coefficient-sign realization \(m_\xi\) of the finite pair-cosine prior in \cref{obj:def:cosine-prior}, with the cutoff and normalization displayed there, is centered and belongs to \(\mathcal H_{d,s}\). Drawing its coefficient signs independently of the covariate sample gives, for every admissible design \(\pi\),
\[
c_sa_s(n,d)
\leq
\mathbb E_\xi\!\left[\mathcal L_{n,d,s}(\pi,m_\xi)\right].
\]

For every admissible \(\pi\) and every \(m\in\mathcal H_{d,s}\), the signing loss is finite, and the actual Horvitz--Thompson estimator under the Gaussian completion satisfies
\[
n\operatorname{Var}(\widehat\theta)-4.01
=
\mathcal L_{n,d,s}(\pi,m).
\]
Consequently, taking the infimum over admissible designs and the supremum over Gaussian completions indexed by \(m\in\mathcal H_{d,s}\) gives exactly \(R_s(n,d)\); taking only the supremum for \(\pi^\star_{n,d}\) gives exactly its worst signing loss.

Let \(\mathcal Q_{d,s}\) be the frozen uniform outcome-law class: the covariate is uniform on \([0,1]^d\), both potential outcomes have finite second moments, and their conditional midpoint mean is a centered member of \(\mathcal H_{d,s}\). For a law \(P\) in this class, write
\[
m_P(u)=\mathbb E_P[(O(1)+O(0))/2\mid U=u],
\qquad
\widehat T_P=\frac2n\sum_{i=1}^n Z_iO_i((Z_i+1)/2),
\]
and
\[
\mathsf V(P)
=
\operatorname{Var}_P\!\left(
\mathbb E_P[O(1)-O(0)\mid U]\right)
+
2\mathbb E_P\operatorname{Var}_P(O(1)\mid U)
+
2\mathbb E_P\operatorname{Var}_P(O(0)\mid U).
\]
Here the estimator uses independent original-unit copies under \(P\), with assignment generated by \(\pi\). The exact published-law minimax excess is
\[
\inf_{\pi\in\mathcal D_{n,d,s}}
\sup_{P\in\mathcal Q_{d,s}}
\max\!\left\{
n\operatorname{Var}_{P,\pi}(\widehat T_P)-\mathsf V(P),\,0
\right\}
=
R_s(n,d).
\]

Finally, for every fixed \(0<s\leq1\) and every integer dimension sequence \((d_n)\) satisfying \(d_n\geq2\) for every \(n\),
\[
R_s(n,d_n)\longrightarrow0
\quad\Longleftrightarrow\quad
\frac{\binom{d_n}{2}}{n}\longrightarrow0
\quad\Longleftrightarrow\quad
d_n=o(\sqrt n),
\qquad n\longrightarrow\infty.
\]
\end{theoremv}
% CAUSALSMITH-CITED-SCOPE-BEGIN thm:capacity-answer
\verificationfootnotetext{\textbf{Formalization scope.} This result uses the cited conclusion \citep{BartlettMendelson2002Complexities} (Theorem 12(4), with Rademacher complexity as in Definition 2) and \citep{Cytrynbaum2026Limits} (Proposition 2.3 (Excess Variance), equation (2.4), with Assumption 2.1 and Definition 2.2). The cited source proof is not formalized here: Lean verifies its use and all remaining steps. If that cited conclusion is false, the portions of this result that depend on it are not certified.}
% CAUSALSMITH-CITED-SCOPE-END thm:capacity-answer

The common scale expresses the number of pairwise interactions relative to the available units. When that ratio is small, the spectral assignment controls prognostic imbalance at a rate determined by smoothness; the independent coefficient prior gives a matching comparison over every admissible design. The exact variance identities then translate this balance criterion into causal precision. The same finite assignment procedure supplies the upper comparison throughout the stated smoothness range.

\subsection{The common comparison scale}

The spectral assignment bound combines the pooled reflection budget in \cref{obj:lem:exact-reflection-budget}, the imbalance identity in \cref{obj:lem:reflection-transfer}, and the assignment guarantees in \cref{obj:lem:ordered-prefix-rounding}. The independent coefficient prior supplies the full-design converse in \cref{obj:thm:full-design-lower}. Both comparisons use the same saturated scale identified in \cref{obj:thm:capacity-answer}. Here the function and design classes are those of \cref{obj:def:sobolev-class,obj:def:design-class}, and the loss is specified in \cref{obj:thm:log-free-bivariate-capacity}. The assignment procedure in \cref{obj:def:capacity-handle} attains this comparison simultaneously across the stated smoothness range.

For every fixed \(s\in(0,1]\), the dimension conclusion of \cref{obj:thm:capacity-answer} applies to every integer sequence \(d_n\geq2\): vanishing minimax scaled excess is equivalent to a vanishing interaction-to-sample ratio and to \(d_n=o(\sqrt n)\). Thus smoothness determines the rate within the subcritical regime, while the threshold for vanishing minimax scaled excess is expressed by the interaction count relative to sample size. The equivalence accommodates arbitrary dimension sequences under fixed smoothness.

\subsection{Outcome-law capacity}

The causal interpretation rests on the exact transfer in \cref{obj:lem:published-prognostic-capacity}. For every admissible design, the uniform outcome-law class \(\mathcal Q_{d,s}\) defined there satisfies
\[
\sup_{P\in\mathcal Q_{d,s}}\mathcal V_n(\pi,P)
=
\sup_{m\in\mathcal H_{d,s}}\mathcal L_{n,d,s}(\pi,m).
\]
The outcome-law criterion therefore has the same minimax value \(R_s(n,d)\), and its worst-case value for \(\pi^\star_{n,d}\) equals that procedure's worst signing loss. These equalities retain finite potential-outcome second moments, uniform covariates, centered prognosis, and the assignment conditions of \cref{obj:def:design-class}. The transfer uses the excess-variance identity of \citet[Proposition 2.3, equation (2.4), under Assumption 2.1 and Definition 2.2]{Cytrynbaum2026Limits}.

The Gaussian completions in \cref{obj:synth_4} realize the entire prognostic class. As recorded in \cref{obj:thm:capacity-answer}, restricting the outcome-law supremum to those completions preserves both the minimax capacity and the worst-case loss of the specified procedure. For the larger bounded-density class defined in \cref{obj:lem:published-prognostic-capacity}, the class transfer instead supplies a lower comparison. With fixed density bounds \(0<\lambda\leq1\leq\Lambda<\infty\) and fixed smoothness, vanishing minimax scaled excess over that class requires \(d_n=o(\sqrt n)\).

\subsection{Relative-excess precision requirements}

Under the Gaussian completion, \cref{obj:lem:ht-transfer} identifies signing loss with actual scaled HT variance above \(V^*=4.01\). Consequently, the upper comparison in \cref{obj:thm:capacity-answer} gives the following sufficient precision requirement, for any \(\varepsilon>0\):
\[
n\geq B\max\left\{
1,\left(\frac{3072}{4.01\varepsilon}\right)^{1/s}
\right\}.
\]
At this sample size, the procedure \(\pi^\star_{n,d}\) satisfies
\[
\sup_{m\in\mathcal H_{d,s}}
n\operatorname{Var}(\widehat\theta)
\leq4.01(1+\varepsilon),
\]
with each variance evaluated under the corresponding Gaussian completion.

The lower comparison in \cref{obj:thm:capacity-answer} gives the necessary requirement for \(0<\varepsilon<c_s/4.01\):
\[
R_s(n,d)\leq4.01\varepsilon
\quad\Longrightarrow\quad
n\geq B\left(\frac{c_s}{4.01\varepsilon}\right)^{1/s}.
\]
It therefore applies to any admissible design attaining the stated relative-excess target uniformly over the prognostic class. Within this target range, the sufficient and necessary requirements both scale with the coordinate-pair count and with \(\varepsilon^{-1/s}\), up to the displayed constants.

\paragraph{Verification note.}
The theorem statements and proofs are machine-checked in Lean 4 under their stated hypotheses. Uniform sampling, exact order-two prognosis, the pooled Sobolev budget, conditional fairness, and outcome-independent assignment are stipulated model conditions in \cref{obj:ass:uniform-draw,obj:ass:order-two,obj:ass:pooled-budget,obj:ass:fair,obj:ass:assignment-independence}; the capacity bounds, transfer equalities, assignment guarantees, and precision implications are proved conclusions. The theorem-local verification disclosures record the cited Rademacher contraction conclusion \citep[Theorem 12(4), with Rademacher complexity as in Definition 2]{BartlettMendelson2002Complexities} and the cited excess-variance identity \citep[Proposition 2.3, equation (2.4), under Assumption 2.1 and Definition 2.2]{Cytrynbaum2026Limits}, wherever each is used. Those published conclusions enter as premises whose source proofs were not formalized: Lean checks their applications and the remaining deductions, with certification of dependent conclusions conditional on the cited premises.
